\bookchapter{제16장: 계정 체계와 클라우드 서비스}

\chapterintro{제 친구 중 한 명이 작년에 스마트폰을 도난당했습니다. 그가 가장 먼저 생각한 것은 “스마트폰이 얼마짜리인가”가 아니라 “안에 있는 사진을 되찾을 수 있을까”였습니다. 안타깝게도 그는 클라우드 백업을 켜지 않았습니다. 스마트폰은 외부 물건이지만, 계정 안의 데이터는 정말 잃어버리면 다시는 돌아오지 않는 것입니다. 이 장에서는 당신의 디지털 라이프에 진짜 보험을 들어보겠습니다.}

\booksection{하나, 당신의 스마트폰 계정은 도대체 얼마나 많은 것을 관리하나요?}

많은 사람이 계정은 앱을 다운로드하는 용도라고 생각합니다. 사실, 스마트폰 계정 하나에 연결된 것은 상상보다 훨씬 많습니다:

\begin{itemize}
  \item \textbf{연락처}: 몇 년 심지어 십수 년간 저장한 전화번호.
  \item \textbf{사진과 비디오}: 아이 탄생부터 졸업식까지, 첫 데이트부터 신혼여행까지.
  \item \textbf{메모}: 즉흥적으로 적은 영감, 일기, 계정 비밀번호(맞습니다, 많은 사람이 비밀번호를 메모에 적어둡니다).
  \item \textbf{앱 구매 기록}: 수백 위안에 산 유료 앱과 인앱 구매.
  \item \textbf{스마트폰 찾기}: 스마트폰 분실 후, 위치 확인, 잠금, 원격 데이터 삭제를 도와주는 유일한 도구.
\end{itemize}

Apple 스마트폰을 사용한다면, 위 모든 것은 Apple ID에 바인딩됩니다. Android를 사용한다면, 화웨이/샤오미/OPPO/vivo 계정과 Google 계정에 바인딩됩니다. 이 계정들은 당신 디지털 라이프의 “호구부(주민등록)”이며, 호구부를 잃으면 많은 것을 설명하기 어렵습니다.

\booksection{둘, 반드시 해야 할 세 가지, 지금 바로}

\begin{enumerate}
  \item \textbf{2단계 인증 활성화}: Apple이든 Android든, 첫 번째는 계정 설정에서 "2단계 인증" 또는 "2단계 검증"을 켜는 것. 이것은 문에 잠금장치를 추가하는 것과 같습니다——비밀번호만으로는 부족하고, 스마트폰으로 인증 코드를 받아야 함. 누군가 인터넷에서 당신의 비밀번호를 추측해도, 스마트폰이 없으면 들어올 수 없습니다.
  \item \textbf{신뢰할 수 있는 복구 방식 연결}: 2단계 인증은 안전하지만, 위험도 있습니다——당신의 스마트폰을 잃어버리면 어떡하죠? 그래서 시스템은 "복구 키" 또는 "예비 이메일/전화번호" 설정을 안내합니다. 이것은 긴급 상황에서 계정을 되찾을 유일한 경로이므로, 가능하면 종이에 적어 스마트폰에 저장하지 마세요.
  \item \textbf{"스마트폰 찾기" 활성화 확인}: 이 기능은 스마트폰 분실 후 위치 확인, 잠금, 소리 재생, 심지어 원격 데이터 삭제를 도와줍니다. 마지막 방어선입니다. 많은 사람이 스마트폰 구매 시 이 단계를 건너뛰고, 나중에 필요할 때 켜지 않은 것을 발견하면 너무 늦습니다.
\end{enumerate}

\booksection{셋, iCloud: Apple 사용자의 “공기”}

Apple 스마트폰을 사용하는 사람은 언젠가 한 가지 영적 고문에 직면합니다: iCloud 공간이 가득 차면 어떡하지?

Apple은 무료로 5GB를 주지만, 솔직히 이 5GB는 한 번 백업하기에도 부족합니다. 그래서 대부분 결국 그 성가신 팝업을 받습니다: “iCloud 저장 공간이 가득 찼습니다”. 이때 세 가지 선택이 있습니다:

\begin{itemize}
  \item \textbf{정리}: 설정 본인 이름 iCloud 저장 공간 관리에서 누가 공간을 차지하는지 확인. 보통 사진과 백업이 대부분. 불필요한 구 백업부터 삭제.
  \item \textbf{업그레이드}: 50GB 월 6위안, 200GB 월 21위안. 개인적으로 50GB가 대부분에게 충분하다고 봅니다, 주로 백업과 사진 썸네일 저장용.
  \item \textbf{우회 방안}: 돈을 쓰고 싶지 않다면, 사진을 정기적으로 컴퓨터로 내보내고 iCloud 사진 동기화를 끄세요. 하지만 그러면 사진 자동 백업의 편리함을 잃게 됩니다.
\end{itemize}

중요한 알림: iCloud 사진 라이브러리는 “네트워크 드라이브”가 아니라——“동기화”입니다. 스마트폰에서 사진을 삭제하면, iCloud에서도 삭제됩니다. “클라우드에 백업하고 스마트폰을 비우고 싶다”면 iCloud 사진은 당신이 원하는 게 아닙니다, iCloud 클라우드 백업이나 서드파티 네트워크 드라이브가 필요합니다.

\begin{tcbinfo}{iCloud+: 유료 사용자만 누릴 수 있는 프라이버시 3종 세트}
임의 유료 iCloud 등급(50GB/200GB/2TB 등)으로 업그레이드하면 자동으로\textbf{iCloud+}사용자가 되어, 세 가지 프라이버시 보호 기능이 잠금 해제됩니다:

\begin{itemize}
  \item \textbf{Private Relay(비밀 릴레이)}: Apple판 "VPN". Safari 트래픽이 Apple 서버를 거쳐 중계된 후 웹사이트로 전송, 누구도(Apple 포함) 당신이 누군지, 무엇을 방문했는지 동시에 볼 수 없습니다. 공공 Wi-Fi에서 프라이버시 보호에 적합. 설정 경로:\texttt{설정 > 본인 이름 > iCloud > 비밀 릴레이}.
  \item \textbf{Hide My Email(이메일 숨기기)}: 웹사이트/앱 가입 시 일회용 무작위 이메일 주소 생성, 메일은 자동으로 진짜 이메일로 전달. 더 이상 이 웹사이트의 메일을 받고 싶지 않으면 이 무작위 주소를 바로 끄면 됨, 진짜 이메일 노출 안 됨. 설정 경로:\texttt{설정 > 본인 이름 > iCloud > 이메일 숨기기}.
  \item \textbf{사용자 지정 이메일 도메인}: 본인 도메인(예:@yourname.com)으로 iCloud 메일 송수신 가능, 개인 브랜드나 가족 이메일 관리에 적합.
\end{itemize}

알림: Private Relay는 일부 국가/지역에서 사용 불가, 국내 사용자는 다른 방안 보완 필요할 수 있음.

\end{tcbinfo}

\begin{tcbwarn}{iCloud 고급 데이터 보호(ADP)——진짜 엔드투엔드 암호화}
2022년 말 Apple은\textbf{iCloud 고급 데이터 보호(Advanced Data Protection, ADP)}를 출시, 엔드투엔드 암호화(E2EE)를 기존 건강, 비밀번호 등 소수 데이터에서\textbf{iCloud 데이터 대부분}으로 확장——사진, 메모, iCloud 클라우드 백업, 기기 정보 등 총 23류 데이터.

\begin{itemize}
  \item \textbf{E2EE란}: 암호화 키는 본인 기기에만 존재, Apple 서버에도 없음. Apple이 해킹당하거나 법 집행 기관이 제출을 요구해도, Apple은 해독 가능한 평문 데이터를 내놓을 수 없음.
  \item \textbf{활성화 경로}:\texttt{설정 > 본인 이름 > iCloud > 고급 데이터 보호}, 안내에 따라 복구 연락처/복구 키 설정(\textbf{필수}, 키 잃어버리면 데이터도 되찾을 수 없음).
  \item \textbf{주의 사항}: (1)활성화 후 반드시 복구 키를 잘 보관, 그렇지 않으면 계정 분실 시 데이터 되찾기 불가; (2)공유 앨범, iCloud 공유 사진 라이브러리 등 소수 기능은 E2EE에 참여하지 않음; (3)중국 본토 Apple ID는 아직 ADP 미지원, 다른 지역 계정으로 전환 필요.
\end{itemize}

한마디 조언: 계정에 중요한 프라이버시 데이터가 있고, 계정 지역이 지원한다면 ADP 활성화 강력 권장——일반인이 얻을 수 있는 가장 강력한 클라우드 데이터 보호입니다.

\end{tcbwarn}

\booksection{넷, 동기화와 백업, 헷갈리지 마세요}

이 두 단어는 비슷하게 들리지만, 결과는 완전히 다릅니다:

\begin{itemize}
  \item \textbf{동기화}는 양방향: 스마트폰에서 삭제하면 클라우드에서도 삭제; 클라우드에서 삭제하면 스마트폰에서도 삭제. iCloud 사진, 연락처, 메모는 모두 동기화. 장점은 다중 기기 간 일관성 유지, 단점은 "삭제하면 진짜 사라짐".
  \item \textbf{백업}은 단방향: 스마트폰이 현재 상태의 "스냅샷"을 클라우드에 저장. 이후 스마트폰에서 무언가를 삭제해도, 클라우드 백업의 사본은 남음. iCloud 클라우드 백업, 컴퓨터 로컬 백업이 모두 백업. 장점은 데이터가 더 안전, 단점은 공간 차지, 그리고 실시간 업데이트 불가.
\end{itemize}

\booksection{다섯, Android 사용자의 계정 매트릭스}

Android 쪽 상황은 Apple보다 복잡합니다, 데이터가 여러 계정에 분산될 수 있기 때문:

\begin{itemize}
  \item \textbf{브랜드 계정}(화웨이/샤오미/OPPO/vivo): 클라우드 공간, 스마트폰 찾기, 앱 스토어, 지갑 관리. 각 브랜드마다 자체 클라우드 서비스, 보통 기본 무료+확장 유료 모델.
  \item \textbf{Google 계정}: Gmail, Google 포토, Google 연락처를 사용한다면, 이것들은 모두 Google 계정에 있습니다. Google은 무료 15GB 제공, Apple보다 관대하지만, Gmail, Drive와 공유.
  \item \textbf{위챗/QQ}: 이들은 스마트폰 시스템과 독립적. 채팅 기록과 위챗 즐겨찾기는 브랜드 계정에도 Google 계정에도 없고, 위챗 자체 서버에만 존재.
\end{itemize}

그래서 Android 사용자가 백업할 때는: 브랜드 클라우드 백업, Google 백업, 위챗 채팅 기록 백업 세 가지를 각각 확인해야, 하나라도 빠지면 안 됩니다.

\booksection{여섯, 가족 공유: 가족도 혜택 보게}

Apple의 “가족 공유”는 가족과 앱 구매, 음악, 영화, iCloud 공간을 공유할 수 있고, 위치도 공유 가능. 최대 6명 추가. 설정 방법: 설정 본인 이름 가족 공유. 화웨이와 샤오미에도 유사 기능, 계정 설정에서 “가족 공유” 또는 “가족 클라우드”를 찾을 수 있음.

매우 실용적인 시나리오: 부모님 스마트폰에 가족 공유를 설정하면, 언제든 그들의 위치를 볼 수 있어(그들이 동의한다는 전제하에), 노인이 길을 잃거나 스마트폰이 연락 안 될 때 위치 확인 가능.

\begin{tcbwarn}{계정 도난의 결과는 상상보다 심각}
해커가 계정 비밀번호를 손에 넣으면, 다음을 할 수 있습니다: 원격에서 스마트폰 잠금 후 몸값 요구, 사진과 연락처 열람, 결제 수단으로 물건 구매. 더 골치 아픈 건, 계정에 너무 많은 것을 연결했다면, 계정 자체를 되찾는 것도 악몽. 그러니 2단계 인증, 지금 바로 켜세요.

\end{tcbwarn}

\begin{tcbpractice}{실습: 계정 종합 건강 검진}
\begin{enumerate}
  \item 스마트폰 설정 열기, 계정 페이지 진입, 2단계 인증 활성화 확인.
  \item "신뢰할 수 있는 전화번호"와 "복구 키"가 설정되어 있는지 확인, 키는 종이에 적기.
  \item iCloud 또는 클라우드 공간 관리 진입, 어떤 데이터가 동기화 중인지 확인, 불필요한 것은 끄기.
  \item "내 스마트폰 찾기"가 활성화되었는지 확인.
  \item Android를 사용한다면, 브랜드 클라우드 백업, Google 백업, 위챗 백업 세 가지 모두 커버되었는지 확인.
\end{enumerate}

\end{tcbpractice}
